% Currículo completo — Português (Brasil). Base: trajetória completa + inputs de IA e founder.
\documentclass[11pt,a4paper]{article}
\usepackage[T1]{fontenc}
\usepackage[utf8]{inputenc}
\usepackage[brazil]{babel}
\usepackage[margin=1.4cm]{geometry}
\usepackage{helvet}
\renewcommand{\familydefault}{\sfdefault}
\usepackage{enumitem}
\usepackage{titlesec}
\usepackage{xcolor}
\usepackage[hidelinks]{hyperref}

\definecolor{ink}{HTML}{111111}
\definecolor{muted}{HTML}{555555}
\color{ink}
\pagestyle{empty}
\setlist[itemize]{leftmargin=1.1em,itemsep=1pt,topsep=2pt}

\titleformat{\section}{\large\bfseries\scshape}{}{0em}{}[{\color{ink}\titlerule[1.2pt]}]
\titlespacing{\section}{0pt}{1.0em}{0.45em}

\newcommand{\job}[3]{\vspace{3pt}\noindent{\bfseries #1} \textcolor{muted}{— #2}\hfill\textcolor{muted}{\small #3}\par}
\newcommand{\heads}[1]{\textbf{#1:} }

\begin{document}

{\Huge\bfseries Danrley Pereira}\\[2pt]
{\large\color{muted} Fundador e Engenheiro de Software Full-Stack — IA, Backend e Produto}\\[4pt]
{\small
danrley.pereira@dwcorp.com.br \textbullet{} +55 61 9 9423-4712 \textbullet{} Brasília, Brasil \textbullet{}
\href{https://www.linkedin.com/in/danrleypereira}{LinkedIn} \textbullet{}
\href{https://github.com/danrleypereira}{GitHub} \textbullet{}
\href{https://danrleypereira.com}{danrleypereira.com}
}

\section{Resumo}
Engenheiro de software full-stack desde 2017 e fundador da DW Corp, construindo \textbf{produtos habilitados por IA de ponta a ponta} — agentes LLM, pipelines RAG e machine learning — sobre uma base sólida de backend em Python, C\#/.NET e Node.js. Entreguei recursos de IA para clientes reais, construí \textbf{fundações de produto reutilizáveis} (autenticação, pagamentos, CI/CD, infraestrutura em nuvem e on-premise), contratei e mentorei engenheiros, e lidero a primeira comunidade de \textbf{Software Craftsmanship} de Brasília. Sou responsável por todo o caminho, dos dados e da recuperação aos serviços de backend e às interfaces web que os entregam, e tenho ampla experiência em software operacional e de logística.

\section{Principais Habilidades}
\begin{itemize}
\item \heads{IA / ML} APIs de LLM (OpenAI, Anthropic, Azure OpenAI), RAG e embeddings, FAISS, LangChain / LangGraph, fluxos com agentes, prompt engineering, TensorFlow / Keras, aprendizado por reforço.
\item \heads{Backend} Python (FastAPI, Flask, Django, Celery), C\#/.NET e ASP.NET, Node.js, Rust, APIs RESTful, arquitetura orientada a eventos (RabbitMQ), microsserviços.
\item \heads{Frontend} React, Next.js, Vue / Nuxt, Angular, TypeScript.
\item \heads{Dados} PostgreSQL, SQLAlchemy, MySQL, MongoDB, Redis, Snowflake, ETL / pipelines de dados, Pandas.
\item \heads{Nuvem e DevOps} Docker, Kubernetes, AWS, GCP, Azure, CI/CD (GitHub e GitLab), Terraform, Grafana, on-premise / Linux bare-metal.
\item \heads{Testes} Pytest, Playwright, Cypress, Cucumber, Behave, Allure, TDD / BDD.
\end{itemize}

\section{Experiência}

\job{Fundador e Engenheiro de IA}{DW Corp}{set 2022 -- Atual}
\begin{itemize}
\item Construí \textbf{agentes LangChain / LangGraph} — um assistente de orçamento no WhatsApp e um barramento multiagente orientado a eventos — atrás de uma interface comum multi-provedor de LLM.
\item Construí o \textbf{Curupira}, uma plataforma de autenticação Rust/OIDC reutilizável (multi-tenancy, SSO, MFA), além de integrações de pagamento (Stripe, Adyen) e um stack bare-metal + nuvem híbrida com Terraform e Grafana.
\item Entreguei produtos para clientes: um gestor de pedidos de restaurantes (Django, Celery, resiliência com circuit-breaker para um data lake) e um app de field-service em .NET MAUI com procedures, views e triggers de banco.
\item Contratei e mentorei mais de 10 engenheiros, e lidero a primeira comunidade de Software Craftsmanship de Brasília (katas de TDD, mesas-redondas).
\end{itemize}

\job{Engenheiro de Software Sênior}{Epicor}{nov 2023 -- mar 2025}
\begin{itemize}
\item Integrei o \textbf{Azure OpenAI} e um \textbf{pipeline RAG} sobre dados proprietários; usei o Azure Document Intelligence para extração estruturada.
\item Construí microsserviços \textbf{.NET Aspire} com processamento orientado a eventos (\textbf{RabbitMQ}), OAuth2/SAML personalizado, implantados em Kubernetes.
\item Liderei a refatoração de um portal Angular para microfrontends, com forte cobertura xUnit e end-to-end, e conduzi o planejamento de sprints e a documentação.
\end{itemize}

\job{Pesquisador}{Ibict / MCTI (governo federal)}{mar 2025 -- Atual}
\begin{itemize}
\item Construí pipelines de dados no \textbf{Snowflake} sobre dados institucionais, avançando rumo ao treinamento de modelos com dados proprietários.
\item Coorientei um TCC que projetou um \textbf{sistema RAG} para o judiciário brasileiro (nota 9,0/10); auditei uma plataforma com mais de 600 mil usuários frente à ISO/IEC 25000 e LGPD.
\end{itemize}

\job{Consultor Full-Stack Sênior}{SubSurface Dynamics}{jan 2024 -- jun 2024}
\begin{itemize}
\item Refatorei um app MERN para um projeto \textbf{TypeScript} moderno sobre \textbf{Fastify, Socket.io, Redis} e Typegoose, atendendo milhares de usuários em tempo real.
\item Ajustei o Redis e as séries temporais do MongoDB Atlas por custo e desempenho; protegi a plataforma com Azure AD (ADAL) + Auth0 (OAuth2/SAML).
\item Construí gráficos 2D/3D e uma visualização de poços com \textbf{Three.js}, além de modo escuro e um sistema de notificações.
\end{itemize}

\job{Engenheiro de Dados e DevOps}{Dado Global}{dez 2021 -- abr 2023}
\begin{itemize}
\item Projetei um sistema \textbf{ETL} com responsabilidades em camadas usando \textbf{SQLAlchemy} e cobertura rigorosa de pytest; automatizei a extração com Selenium e Behave.
\item Levei a entrega a \textbf{Continuous Delivery} com GitHub Actions (semanal $\rightarrow$ diária), melhorando o SLA de 70\% para 95\%.
\end{itemize}

\job{Engenheiro de Software}{ClearSale}{nov 2021 -- mar 2022}
\begin{itemize}
\item Construí microsserviços antifraude de alto desempenho em \textbf{C\#/.NET}; projetei views, triggers e stored procedures e otimizei consultas com Dapper.
\item Defini processos de engenharia (UML/BPMN), iniciei testes unitários e mentorei o time.
\end{itemize}

\job{Líder Técnico e Pesquisador}{LabTech, UDF}{jan 2023 -- dez 2023}
\begin{itemize}
\item Liderei uma plataforma multicamada de gestão de eventos (React, Python, Spring Boot, RabbitMQ, Kubernetes) e mentorei engenheiros juniores.
\end{itemize}

\job{Engenheiro de Software}{Layers Education}{mar 2022 -- set 2022}
\begin{itemize}
\item Mantive sistemas backend e microfrontend; reduzi os tempos de consulta do MongoDB de 1,3s para 0,6s sob carga, construí um data lake no BigQuery para ETL e reforcei a segurança conforme OWASP.
\end{itemize}

\job{Engenheiro Full-Stack}{Compuletra}{jul 2020 -- ago 2021}
\begin{itemize}
\item Construí software de logística de \textbf{transporte, entrega e pátio} com sistemas de frota e \textbf{inspeção de veículos}, e uma plataforma ERP desktop (Electron) sobre PostgreSQL.
\item Migrei um app legado de rastreamento veicular de AngularJS para React (on-premise), e construí um gerador de páginas comerciais com checkout compatível com PCI gerenciado a partir do banco.
\end{itemize}

\job{Desenvolvedor Full-Stack e BPM}{ConsultMídia / Banco de Brasília}{abr 2020 -- jan 2021}
\begin{itemize}
\item Construí interfaces AngularJS/Spring Boot e APIs REST sobre BPM \textbf{Camunda} para modernizar os fluxos do banco; integrei o Alfresco (gestão documental) e automatizei testes de desempenho (JMeter, Selenium).
\end{itemize}

\job{Desenvolvedor Full-Stack e Mobile}{Transoft}{fev 2021 -- out 2021}
\begin{itemize}
\item Construí e renovei aplicações para gerenciar sensores e câmeras de frotas de ônibus, e um app mobile de gestão de pessoal integrado ao Oracle (Angular, Ionic, Laravel).
\end{itemize}

\job{Desenvolvedor de Plataforma}{Ministério da Saúde / USP}{jun 2019 -- jul 2020}
\begin{itemize}
\item Ajudei o Ministério da Saúde a gerir a \textbf{distribuição de medicamentos em mais de 5.000 municípios}; construí o frontend e integrei APIs REST e microsserviços para expor dados, reduzir fraudes e otimizar a integração de serviços.
\end{itemize}

\section{Experiência Anterior}
\begin{itemize}
\item \textbf{Taurus Software} — aplicações de gestão de energia IoT, mobile híbrido (Ionic) e de captura de leads. \emph{2018--2020}
\item \textbf{Gama Cidadão} — Gerente Web e de Marketing Digital; portais WordPress/Vue, landing pages e SEO. \emph{2017--2018}
\item \textbf{UMESB} — Desenvolvedor por contrato (sistema de matrícula de estudantes, Polymer/Node.js). \emph{2017}
\item \textbf{IT Universe} — Instrutor de web/design gráfico e fundamentos de computação. \emph{2017}
\end{itemize}

\section{Formação}
\begin{itemize}
\item Bacharelado em Engenharia de Software — Universidade de Brasília
\item Bacharelado em Ciência da Computação — Centro Universitário do Distrito Federal
\item MBA, Cadeia de Suprimentos e Logística Integrada — Southern Cross University
\item MBA, Administração de Banco de Dados — Anhanguera
\end{itemize}

\section{Certificações e Palestras}
\begin{itemize}
\item Palestrante no palco principal da \textbf{Campus Party Brasil} — Clean Architecture, TDD, IA, SDET e Software Craftsmanship.
\item Business Intelligence com Power BI \textbullet{} Scrum Fundamentals (SFC) \textbullet{} Kanban Foundation \textbullet{} Produção Oral B2 (Inglês) — UnB.
\end{itemize}

\section{Idiomas}
Português (nativo) \textbullet{} Inglês (C2) \textbullet{} Espanhol (intermediário)

\end{document}
